\chapter{One Cosine Integral, Two Checked Proofs}
\label{ch:two-cosine-proofs}

\begin{center}
\href{proof-comparison/index.html}{\textbf{Open the measured proof comparison and interactive dependency graph}}
\end{center}

This chapter records two named Lean proofs of exactly the same proposition.
The computation is fixed independently of either proof.  The interactive
comparison extracts references from stored declaration types and bodies;
it is not an import graph.  The manuscript graph below is a selected
mathematical overview, while every edge in the interactive overview has
an inspectable path through actual Lean references.

\begin{definition}[The common cosine-integral statement]
\label{def:cosine-comparison-statement}
Let $B$ be the existing certified inverse-arctangent provider.  On rational
$0\le a\le b\le1/2$, let $S,C$ be its geometric sine and cosine, with
normalized angle.  Let $R_{k,n}$ be the existing cosine left rectangle sum
on $k+1$ equal cells, evaluated at stage $n$, and set
\[
 I_n=\bigcap_{k=0}^{n}\operatorname{expand}
       \left(R_{k,n},\frac{4000(b-a)^2}{k+1}\right).
\]
The common proposition is the equivalence of this integral program with
$\operatorname{reciprocalPiRaw}(S(b)-S(a))$.
Its Lean name is \code{CosineFTC.IntegralCosineStatement}.
No integral endpoint, concavity, or derivative certificate is an additional
hypothesis beyond the data already needed to define $S,C$.
\end{definition}

\begin{lemma}[The shared arctangent clock]
\label{lem:cosine-comparison-clock}
\leanok
The rectangle computation of $\int_0^u(1+t^2)^{-1}\,dt$ is equivalent to
the geometric arctangent computation.  The inverse provider computes the
half-angle parameter satisfying $A(u)\simeq\pi x/2$.  Rational circle
coordinates then compute $S,C$.  This foundation is shared by both proofs.
\end{lemma}

\begin{lemma}[Finite quadratic residual]
\label{lem:cosine-comparison-residual}
\uses{lem:cosine-comparison-clock}
\leanok
For a fixed positive rational increment $h$ in the chart, sufficiently
late rational selections from the sine, cosine, and pi boxes satisfy
$|s_1-s_0-hpc|\le4000h^2$.  This follows from rational circle and clock
inequalities without invoking the normalized sine derivative or concave FTC.
\end{lemma}

\begin{theorem}[First proof: direct inequalities]
\label{thm:cosine-comparison-direct}
\uses{def:cosine-comparison-statement,lem:cosine-comparison-residual}
\leanok
The common cosine-integral statement holds.  The proof
\code{CosineFTC.integral\_cosPi\_viaInequalities} telescopes finitely many
rational sine increments, uses a shared reciprocal-pi corner, and transports
the resulting bound to every pair of stages by nesting.  Finite
intersections give the equivalence.
\end{theorem}

\begin{lemma}[Geometric supporting secants and concavity]
\label{lem:cosine-comparison-concavity}
\uses{lem:cosine-comparison-clock}
\leanok
Circle chords and tangent slopes give, for $x<y$,
\[
 \pi C(y)\preceq\frac{S(y)-S(x)}{y-x}\preceq\pi C(x).
\]
Removing arbitrary rational slack proves exact concavity of $S$ and
$P=S/\pi$.  This argument uses neither a first nor a second derivative.
The supporting secants for $P$ have slope computation $C$, with the proved
conservative Lipschitz bound $32$.
\end{lemma}

\begin{theorem}[Companion: the constructed secant derivative]
\label{thm:cosine-comparison-derivative}
\uses{lem:cosine-comparison-concavity}
\leanok
For interior rational $x$, let $r=\min(x,1/2-x)$ and $h_k=r/(k+1)$.
Intersect the outward forward/backward secant brackets for $k\le n$,
all evaluated at stage $n$.  Concavity proves compatibility, and the
quantitative slope bounds prove shrinking.  This sine-only derivative
program is valid and equivalent to $\pi C(x)$.
It is a companion consequence of the same certificate: its validity theorem
is not itself a referenced premise of either final integral proof.
\end{theorem}

\begin{theorem}[General concave FTC]
\label{thm:cosine-comparison-general-ftc}
\leanok
Exact concavity, valid supporting-slope computations, cross-stage secant
bounds, and a proved Lipschitz bound suffice to compare finite derivative
rectangle sums with primitive endpoints.  With relevant source widths at
most $\tau$, the local rational estimate is
\[
 |F_n^-(y)-F_n^-(x)-hD_n^-(x)|\le Kh^2+(1+h)\tau.
\]
A common stage for finitely many samples, a finite telescope, and removal
of rational slack establish the integral equivalence.  The certificate
contains no assumed integral overlap or endpoint identity.
\end{theorem}

\begin{theorem}[Second proof: via concave FTC]
\label{thm:cosine-comparison-ftc}
\uses{def:cosine-comparison-statement,lem:cosine-comparison-concavity,thm:cosine-comparison-general-ftc}
\leanok
The identical common statement holds by
\code{CosineFTC.integral\_cosPi\_viaFTC}.  Apply the general concave FTC
to $P=S/\pi$ and $D=C$, then identify the endpoint difference by finite
distributivity.  This proof neither calls the direct proof nor uses its
quadratic-residual theorem.
\end{theorem}

\section{Quantitative comparison}
The \href{proof-comparison/index.html}{interactive comparison} reports the
final proof's stored expression size, the transitively referenced project
declarations, shared and exclusive supporting material, source-range line
counts, and a limited kernel recheck benchmark.  It distinguishes expanded
expression-tree occurrences from structurally distinct subexpressions.
Counts include generated helpers and implicit arguments.  They do not
measure elegance, proof-discovery difficulty, or mathematical strength.

The build verifies identical complete theorem types and the separation of
the two proof paths.  Both share the earlier arctangent/area foundation.
Axiom dependencies remain separately inspectable: absence of admitted
proofs is not a claim that all upstream native computation was replaced by
kernel-only arithmetic.  The web page pins its measured source revision.
